\clearpage\SourceLeaf{529}\renewcommand{\BookRunningTitle}{Bemerkungen.}
\pdfbookmark[0]{Bemerkungen.}{bemerkungen-second-impression}
\begin{center}\textbf{BEMERKUNGEN.}\end{center}
\BookRule{.16\textwidth}
Diesem zweiten Bande von \textsc{Gauss} Werken habe ich alle Abhandlungen, Aufsätze und Tafeln aus dem Gebiete der Höhern Arithmetik, soweit die sieben Sectionen der \textit{Disqu. Arithm.} sie nicht schon umfassen, einverleibt, und zwar die in den ‘\textit{Commentationes societatis regiae scientiarum Gottingensis}’ (in Quart) veröffentlichten fünf Abhandlungen, die in den ‘\textit{Göttingischen Gelehrten Anzeigen}’ (in Octav) erschienenen (von \textsc{Gauss} nicht unterzeichneten, aber durch die Acten der Göttinger Universitäts-Bibliothek in Betreff der Autorschaft verificirten) Anzeigen sowohl dieser eignen als auch einiger anderer nichteigner Schriften, und eine Auswahl aus dem Handschriftlichen Nachlasse.
\par Beim zweiten Abdruck habe ich noch die Tabellen ‘Circuli quadratura nova’ ‘Zur Berechnung der Logarithmen’ ‘Quadratorum myrias prima’ ‘Indices der Primzahlen im höhern Zahlenreiche’ ‘Hülfstafel zur Auflösung der unbestimmten Gleichung $A=fxx+gyy$ vermittelst der Ausschliessungsmethode’ ferner ‘Sectio octava’, so weit sie aufgeschrieben ist und endlich zwei Briefe von \textsc{Gauss} an \textsc{Dirichlet} als wesentliche Stücke der Geschichte der Höhern Arithmetik hinzugefügt.
\par Zur bessern Uebersicht der Gegenstände in einem so umfangreichen Bande sind die Lehrsätze auf gleiche Weise durch den Druck ausgezeichnet. Zum leichtern Gebrauch sowohl der ältern Ausgaben wie der vorliegenden ist bei den Verweisungen auf die Disq. Arithm. statt der Nummer der Seite die der Artikel gesetzt, so wie bei den Angaben von Abhandlungen statt des Orts ihrer Veröffentlichung deren eigner Titel. Die Note, die dem Art. 2 der Abhandlung ‘\textit{Theorematis arithmetici demonstratio nova}’ ursprünglich
